\section*{Chapter \ref{ch:b2:multiple-regression} --- Multiple Regression}

\begin{solution}{exo:b2:multiple-regression:1}
% figdata: steps
4.10 per cent with no control, 5.22 after adding age, 2.04 with every control. The largest single move is the last,
adding log 1975 earnings: from 3.83 to 2.04, 1.79 points; adding age squared moves it 1.23 points, adding age 1.12.
\end{solution}

\begin{solution}{exo:b2:multiple-regression:2}
Men above the cap are recorded at the cap, so high earners appear to earn less than they do. Men with much schooling are
more often above the cap, so the cap compresses the gap between them and men with little schooling: the measured return to
schooling is likely too small.
\end{solution}

\begin{solution}{exo:b2:multiple-regression:3}
% code: exercises/solutions.py adj_r2
$1 - (1 - 0.1181) \times 13{,}054/13{,}048 = 1 - 0.8819 \times 1.00046 = 0.1177$.
\end{solution}

\begin{solution}{exo:b2:multiple-regression:4}
Schooling raises 1975 earnings, and 1975 earnings predict 1978 earnings. Holding 1975 earnings fixed compares men with
different schooling but the same 1975 earnings, which removes the part of the effect of schooling that works through
earlier earnings. A control is bad if it is itself affected by the regressor: evidence would be that the regressor
predicts the control (here a year of schooling goes with 4.5 per cent higher 1975 earnings), or, better, that the control
was measured after the regressor was set.
\end{solution}

\begin{solution}{exo:b2:multiple-regression:5}
% code: exercises/solutions.py X5, Y5, ols, r2, adj_r2
$X'X = \begin{pmatrix}5 & 15 & 3\\ 15 & 55 & 12\\ 3 & 12 & 3\end{pmatrix}$, $X'y = (44, 161, 36)'$. The solution is
$b = (1, 2, 3)'$: check $5 + 30 + 9 = 44$, $15 + 110 + 36 = 161$, $3 + 24 + 9 = 36$. Fitted values $3, 5, 10, 12, 14$;
residuals $1, -1, -1, 1, 0$, so $e'e = 4$. With $\bar y = 8.8$, $\sum(y_i - \bar y)^2 = 90.8$, so $R^2 = 1 - 4/90.8 = 0.956$
and $\bar R^2 = 1 - (4/2)/(90.8/4) = 0.912$.
\end{solution}

\begin{solution}{exo:b2:multiple-regression:6}
% code: tests/test_solutions.py test_age_example
Earnings rise with age ($\gamma > 0$) and older men have less schooling ($\delta < 0$), so leaving age out biases the
schooling coefficient down. The formula: $0.0522 + 0.0192 \times (-0.585) = 0.0522 - 0.0112 = 0.0410$, the short
coefficient.
\end{solution}

\begin{solution}{exo:b2:multiple-regression:7}
% code: exercises/solutions.py brenland, cps
(a) 5.80 per cent per year of schooling with no control, $R^2 = 0.108$; 5.84 per cent with age and age squared,
$R^2 = 0.309$. (b) With \texttt{nodegree} the schooling coefficient falls from 3.83 to 2.16 per cent; men without a
diploma earn 14.8 log points less at the same years of schooling. The standard error rises from 0.0028 to 0.0039, and the
VIF of schooling from 1.14 to 2.24: years of schooling and having a diploma are strongly related.
\end{solution}

\begin{solution}{exo:b2:multiple-regression:8}
% code: exercises/solutions.py collinearity_table
By \eqref{eq:b2:multiple-regression:var} the standard error falls as $1/\sqrt n$: by a factor of 2. The output gives 0.022,
0.025, 0.052 and 0.160 against 0.045, 0.057, 0.103 and 0.311. The ratios are close to 2, though not exactly, because new
draws give new sample correlations.
\end{solution}

\begin{solution}{exo:b2:multiple-regression:9}
The column of ones equals the sum of the two indicators, so the three columns are linearly dependent and the rank is 2.
Either drop one indicator: the constant is then the mean for the omitted group (other regressors fixed) and the remaining
coefficient the difference between the groups. Or drop the constant: each indicator's coefficient is its group's mean.
\end{solution}

\begin{solution}{exo:b2:multiple-regression:10}
Let $P = W(W'W)^{-1}W'$. Then $P' = W((W'W)^{-1})'W' = P$, since $W'W$ is symmetric, and
\[PP = W(W'W)^{-1}W'W(W'W)^{-1}W' = P.\]
So
$M = I - P$ is symmetric and $MM = I - 2P + P = M$. $MW = W - W(W'W)^{-1}W'W = 0$. The regression of $y$ on $W$ has
coefficients $(W'W)^{-1}W'y$ and fitted values $Py$, so its residual is $y - Py = My$.
\end{solution}

\begin{solution}{exo:b2:multiple-regression:11}
Let $Z = (x, W)$ and $A = (Z'Z)^{-1}Z'$, so that the short regression's coefficients are $Ay$. From the long regression,
$y = Z\hat b + z\hat\gamma + e$, where $\hat b$ collects $\hat\beta_1$ and $\hat\beta_W$. Then $Ay = AZ\hat b + Az\,\hat\gamma
+ Ae = \hat b + \hat\gamma Az$, since $AZ = I$ and $Z'e = 0$. The first element of $Az$ is $\hat\delta$, so
$\hat\beta_1^{\text{short}} = \hat\beta_1^{\text{long}} + \hat\gamma\hat\delta$.
\end{solution}

\begin{solution}{exo:b2:multiple-regression:12}
By the theorem $\hat\beta_j = \tilde x_j'y/\tilde x_j'\tilde x_j$ with $\tilde x_j = Mx_j$, a function of $X$ only. With
$y = X\beta + u$, $\tilde x_j'y = \tilde x_j'x_j\beta_j + \tilde x_j'u$ (since $\tilde x_j$ is orthogonal to the other
columns), and $\tilde x_j'x_j = \tilde x_j'\tilde x_j$. So $\hat\beta_j - \beta_j = \tilde x_j'u/\tilde x_j'\tilde x_j$ and
$\operatorname{Var}(\hat\beta_j \mid X) = \sigma^2\tilde x_j'\tilde x_j/(\tilde x_j'\tilde x_j)^2 = \sigma^2/\tilde
x_j'\tilde x_j$. The sum of squared residuals of $x_j$ on the others is $(1 - R_j^2)\sum_i(x_{ij} - \bar x_j)^2 =
(1 - R_j^2)ns_j^2$.
\end{solution}

\begin{solution}{pb:b2:multiple-regression:1}
% code: exercises/solutions.py brenland, brenland_ovb, vif; tests/test_solutions.py test_brenland
\textbf{Part I.}
1. $X$ is $1{,}922 \times 4$; row $i$ is $(1, s_i, a_i, a_i^2)$.
2. $4 \times 4$; $4 \times 1$.
3. The number of heads, 1,922; the sum of $y_i$.
4. The estimator is $(X'X)^{-1}X'y$; if $X'X$ is singular some coefficients are not identified.
5. The residuals sum to zero.
\textbf{Part II.}
6. 0.0580; $R^2 = 0.108$.
7. 0.0584, interval 0.0518 to 0.0649; $R^2 = 0.309$, $\bar R^2 = 0.308$.
8. Among heads of the same age, a year more schooling goes with gross income about 5.8 per cent higher.
9. The coefficients are 0.0724 on $a$ and $-0.000755$ on $a^2$; the peak is at $0.0724/(2 \times 0.000755) = 48$ years.
10. In the survey schooling and age are almost uncorrelated, so leaving age out removes little.
\textbf{Part III.}
11. The variance of $s$ is 6.728, of $\tilde s$ 6.712: age predicts almost none of schooling.
12. 0.0584, the long coefficient, as the theorem says.
13. $\delta_a = 0.318$ and $\delta_{a^2} = 30.99$.
14. $\hat\beta_s^{\text{short}} = \hat\beta_s^{\text{long}} + \hat\gamma_a\hat\delta_a + \hat\gamma_{a^2}\hat\delta_{a^2}$.
15. $0.0724 \times 0.318 - 0.000755 \times 30.99 = 0.0230 - 0.0234 = -0.0004$; $0.0584 - 0.0004 = 0.0580$, the short
coefficient.
16. $a^2$ is measured in squared years, so its coefficients are small and its $\delta$ large; the two terms have opposite
signs and nearly cancel.
\textbf{Part IV.}
17. Experience is an exact linear combination of the constant, $s$ and $a$: $X$ has rank 3 with 4 columns, and the
coefficients are not identified.
18. 3.
19. Only combinations such as $\beta_s - \beta_e$ and $\beta_a + \beta_e$: the effect of schooling at a given age, or of age
at given schooling; never the three separately.
20. 1.00 for $s$, 43.5 for $a$ and for $a^2$.
21. The VIF of $s$ is 1.00: $s$ is almost uncorrelated with $a$ and $a^2$, and its variance does not depend on how well $a$
and $a^2$ can be told apart.
22. By about $\sqrt2 = 1.41$.
23. Upwards: $\gamma > 0$ and $\delta > 0$.
24. No: hours do not measure ability. They may be a bad control, since schooling may change hours worked.
25. The return to a year of schooling in Brenland is 5.84 per cent (95 per cent interval 5.18 to 6.49) holding age fixed;
leaving age out changes it by $-0.04$ points, because age and schooling are almost uncorrelated in the survey.
\end{solution}
